% ===================================================================
%  Minimalbeispiel: Liedblatt – „Alle meine Entchen“ (Volkslied, gemeinfrei)
%  C-Dur, 4/4-Takt. Erste Strophe unter den Noten (musixlyr),
%  Strophen 2–5 zweispaltig darunter, alles auf einer Seite.
%  Melodie: c d e f | g g | a a a a | g | a a a a | g | f f f f | e e | d d d d | c
%  In MusiXTeX ist das a' der Buchstabe h (absolute Tonhöhen, siehe SKILL.md).
%  Eigenständig: braucht keine header.tex.
%
%  Drei Durchläufe:  pdflatex … → musixflx … → pdflatex …
%  (musixflx gleicht alle Notenzeilen auf dieselbe Länge aus.)
% ===================================================================
\documentclass[a4paper,12pt]{article}
\usepackage[T1]{fontenc}
\usepackage[utf8]{inputenc}
\usepackage[ngerman]{babel}
\usepackage{lmodern}
\usepackage[margin=1.8cm]{geometry}
\usepackage{multicol}       % Strophen zweispaltig
\usepackage{musixtex}
\input musixlyr             % Liedtext unter den Noten
\pagestyle{empty}
\setlength{\columnsep}{1em}

% Eine Strophe: \strophe{Nummer}{Text}; bleibt beim Spaltenumbruch zusammen.
\newcommand{\strophe}[2]{%
  \noindent\makebox[1.4em][l]{\textbf{#1.}}%
  \begin{minipage}[t]{\dimexpr\linewidth-1.4em}\raggedright #2\end{minipage}%
  \par\vspace{1.2em}}

\begin{document}
\begin{center}
{\LARGE\bfseries\sffamily Alle meine Entchen}\\[4pt]
{\small Text und Melodie: Volksweise \textbullet\ C-Dur}
\end{center}
\vspace{1em}

\Largemusicsize             % großes System: VOR \begin{music}
\parindent0pt                % erste Zeile nicht einrücken: alle Zeilen gleich lang
\begin{music}
\parindent0pt                % erste Zeile nicht einrücken: alle Zeilen gleich lang
\staffbotmarg=8\Interligne  % Platz unter jedem System (innerhalb von music)
\instrumentnumber{1}\setclef1\treble
\generalmeter{\meterfrac44}\nobarnumbers
\setsongraise1{8mm}         % Text dicht unter die eigene Zeile
\setlyrics{strophe}{Al-le mei-ne Ent-chen schwim-men auf dem See,
schwim-men auf dem See, Köpf-chen in das Was-ser, Schwänz-chen in die Höh.}
\assignlyrics1{strophe}
\startpiece
\Notes\qu c\qu d\qu e\qu f\en\bar      % Al-le mei-ne
\Notes\hu g\hu g\en\bar                % Ent-chen
\Notes\qu h\qu h\qu h\qu h\en\bar      % schwim-men auf dem   (h = a')
\Notes\wh g\en                         % See,
\alaligne
\Notes\qu h\qu h\qu h\qu h\en\bar      % schwim-men auf dem
\Notes\wh g\en\bar                     % See,
\Notes\qu f\qu f\qu f\qu f\en          % Köpf-chen in das
\alaligne
\Notes\hu e\hu e\en\bar                % Was-ser,
\Notes\qu d\qu d\qu d\qu d\en\bar      % Schwänz-chen in die
\Notes\wh c\en                         % Höh.
\Endpiece
\end{music}
\assignlyrics1{}            % Zuweisung beenden (gilt sonst global weiter)

\vspace{2em}

\begin{multicols}{2}
\strophe{2}{Alle meine Täubchen\\ gurren auf dem Dach,\\ gurren auf dem Dach,\\
fliegt eins in die Lüfte,\\ fliegen alle nach.}

\strophe{3}{Alle meine Hühner\\ scharren in dem Stroh,\\ scharren in dem Stroh,\\
finden sie ein Körnchen,\\ sind sie alle froh.}

\strophe{4}{Alle meine Gänschen\\ watscheln durch den Grund,\\ watscheln durch den Grund,\\
suchen in dem Tümpel,\\ werden kugelrund.}

\strophe{5}{Alle meine Püppchen,\\ Emmi und Marie,\\ Emmi und Marie,\\
liegen in der Wiege,\\ bis ich wecke sie.}
\end{multicols}

\end{document}
